\documentclass[10pt,a4paper]{amsart}
\usepackage[margin=19mm]{geometry}
\usepackage{amsmath,amssymb,mathtools}
\usepackage[scr=rsfs,cal=euler]{mathalfa}
\usepackage{xfrac}
\usepackage[unicode,hidelinks,bookmarks=false]{hyperref}
\newcommand{\ie}{i.e.\ }
\newcommand{\cf}{cf.\ }
\newcommand{\resp}{respectively\ }
\newcommand{\Subsection}{\subsection}
\providecommand{\nobreakdash}{\nobreak\hbox{-}}
\setlength{\emergencystretch}{2em}
\multlinegap0pt
\usepackage{bm}
\usepackage{bbm}
\usepackage{dsfont}
\usepackage{xspace}
\usepackage{cancel}
\usepackage{mathtools}
\usepackage{amsthm}
\makeatletter
\@ifundefined{theorem}{\newtheorem{theorem}{Theorem}}{}
\@ifundefined{lemma}{\newtheorem{lemma}[theorem]{Lemma}}{}
\makeatother
\DeclareMathOperator{\supp}{\operatorname{supp}}
\title[Quantitative Polynomial Wiener-Wintner Theorems]{Quantitative Polynomial Wiener-Wintner Theorems}
\author{Lars Becker, Asgar Jamneshan, Christoph Thiele}
\date{Source: arXiv:2601.03999v3 (CC BY 4.0)}
\begin{document}
\maketitle

\noindent\emph{Source and licence.} Quantitative Polynomial Wiener-Wintner Theorems. Version arXiv:2601.03999v3 (CC BY 4.0), distributed under the Creative Commons Attribution 4.0 International licence, as recorded in the arXiv licence metadata for this exact version and shown on the abstract page https://arxiv.org/abs/2601.03999v3. Full rights notice: licenses/arxiv-2601.03999-CC-BY-4.0.txt.\par\smallskip

\noindent\textit{Editorial scope.} Counted unit: the Lemma whose source label is \texttt{l:gamma} in the article (source0.tex lines 545--557, source label \texttt{l:gamma}), delivered together with the article's own complete original proof of it (source0.tex lines 559--609, one \texttt{proof} environment of 51 lines). No mathematics is written, repaired or re-derived: statement and proof are the article's own text line for line. Required context retained uncounted: e:zeta (lines 366--368); e:zeta\_supp (lines 370--372); e:zeta\_int (lines 374--376). Every definition, display and label of those lines is retained unchanged. Omitted from this package and not redistributed: 1--365, 369--369, 373--373, 377--544, 558--558, 610--1559. Nothing inside a retained excerpt is omitted or altered: every retained body is delivered exactly as the article's lines are written, so no omission receipt of any kind accompanies this package. Citations: the retained excerpts carry no citation command at all, so no bibliography entry of the article is transcribed and no reference is redistributed. Reference adaptation: none was needed and none was made; every reference inside the delivered unit resolves inside it, so no unresolved reference or citation is produced. Printed slips kept literal: no printed word, symbol, hypothesis or number of the counted statement or of its proof was altered. Layout and class: the article's own class is \texttt{amsart} with packages including inputenc, xcolor, mathtools, amssymb, amsthm, amsmath, graphicx, enumitem, hyperref; the wrapper is the lane's pinned \texttt{amsart} editorial wrapper at 10pt on A4, which restates the theorem-like environments the retained text needs and adds the article's own macro definitions that the retained text needs (\texttt{\textbackslash supp}), with extra wrapper packages (bm, bbm, dsfont, xspace, cancel, mathtools). The delivered numbering is the unit's own. Assets: no retained span contains a figure environment, a graphics-inclusion command, a PDF-page inclusion, a table or a TikZ picture; no image file is copied into this package, the package declares no asset dependency and no image file is redistributed with it. Licence: version arXiv:2601.03999v3 of this article is distributed under the Creative Commons Attribution 4.0 International licence, as recorded in the arXiv licence metadata for this exact version and shown on the abstract page https://arxiv.org/abs/2601.03999v3. A full rights notice is delivered at licenses/arxiv-2601.03999-CC-BY-4.0.txt. Characters: every retained excerpt is pure ASCII, so the delivered engine, XeLaTeX with its default Latin Modern text font, reports no missing character.
\begin{equation}\label{e:zeta}
    \zeta^\delta(x) = \log(1 + \delta)^{-1} \mathbf{1}_{[1,1 + \delta]}(x). 
\end{equation}

\begin{equation}
    \supp \zeta^{\delta} \subset [1, 1+\delta], \label{e:zeta_supp}
\end{equation}

\begin{equation}
    \int_0^\infty (\zeta^{\delta})_t(x) \, \mathrm{d}t = 1, \qquad x \in (0,\infty).\label{e:zeta_int}
\end{equation}

    \begin{lemma}
        \label{l:gamma}
        It holds that
        \begin{equation}
            \label{e:g_upper}
            |\Gamma_x(s)| \le 1
        \end{equation}
        and, for all $s/2 \le s' \le s$ and $\gamma \le \min\{1/2, \alpha\}$, it holds that
        \begin{equation}
            \label{e:g_hol}
            |\Gamma_x(s) - \Gamma_x(s')| \le 8 \Big(\frac{s - s'}{s'}\Big)^\gamma.
        \end{equation}
    \end{lemma}

    \begin{proof}
    We consider $x, \delta$ fixed and drop them from the notation. Both estimates \eqref{e:g_upper} and \eqref{e:g_hol} scale correctly, so that we may assume $s = 1$ and $s' \in [1/2, 1)$. Since $\|w\|_{\ell^{r'}} \le 1$,
    it holds in particular for all $j$ that
    \[  
        |w_j| \le 1.
    \]
    This along with \eqref{e:zeta}, \eqref{e:zeta_int} implies \eqref{e:g_upper}. To show \eqref{e:g_hol}, we split the indices $j$ into
    \[
        \mathcal{I}_1 = \{j \in \{1, \dotsc, J\} \ : \ u_{j-1} > s' u_{j}\}
    \]
    and 
    \[
        \mathcal{I}_2 = \{j \in \{1, \dotsc, J\} \ : \ u_{j-1}  \le s' u_{j}\}.
    \]
    Using the support assumption \eqref{e:zeta_supp} and $s' \in [1/2,1)$, we have (recall that $\zeta_t = \zeta_t^\delta$)
    \begin{align*}
        &\delta^{-\alpha}|\Gamma(1) - \Gamma(s')|\\
        &= \Big| \int_{1/(2+2\delta)}^{1} \zeta_t(1) \sum_{j \in \mathcal{I}_1 \cup \mathcal{I}_2} w_j \mathbf{1}_{[u_{j-1}, u_j]}(t) - \zeta_t(s') \sum_{j \in \mathcal{I}_1 \cup \mathcal{I}_2} w_j \mathbf{1}_{[u_{j-1}, u_j]}(t)\, \mathrm{d}t \Big|.
    \end{align*}
    We apply the triangle inequality, and a change of variables $t \mapsto \frac{1}{s'}t $ in the $\Gamma(s')$ integral for the $j \in \mathcal{I}_2$ summands to bound this by
    \begin{align}
        &\quad \int_{1/(1+\delta)}^{1} |\zeta_t(1)| \Big|\sum_{j \in \mathcal{I}_1} w_j \mathbf{1}_{[u_{j-1}, u_j]}(t)\Big| \, \mathrm{d}t\label{e:I11}\\
        &+\int_{1/(1+\delta)}^{1} |\zeta_t(1)| \Big|\sum_{j \in \mathcal{I}_1} w_j \mathbf{1}_{[u_{j-1}, u_j]}(\frac{1}{s'}t)\Big| \, \mathrm{d}t\label{e:I12}\\
        &+ \int_{1/(1+\delta)}^{1} |\zeta_t(1)| \Big|\sum_{j \in \mathcal{I}_2} w_j \mathbf{1}_{[u_{j-1}, u_j]}(t) -  \sum_{j \in \mathcal{I}_2} w_j \mathbf{1}_{[u_{j-1}, u_j]}(\frac{1}{s'}t )\Big|\, \mathrm{d}t \label{e:I2}.
    \end{align}
    Using \eqref{e:zeta_int} and Jensen's inequality we bound \eqref{e:I11} by 
    \[
        \Big(\int_{1/(1+\delta)}^1 |\zeta_t(1)| \Big|\sum_{j \in \mathcal{I}_1} w_j \mathbf{1}_{[u_{j-1}, u_j]}(t)\Big|^{\frac{1}{\gamma}} \, \mathrm{d}t \Big)^\gamma.
    \]
    Note that $\zeta_t(1) \le 2 \delta^{-1} t^{-1}$ by \eqref{e:zeta}. This bounds the integral in the previous line by 
    \begin{align*}
        2\delta^{-1} \int \sum_{j \in \mathcal{I}_1} \mathbf{1}_{[u_{j-1}, u_j]}(t) |w_j|^{\frac{1}{\gamma}} \, \frac{\mathrm{d}t}{t} &= \sum_{j \in \mathcal{I}_1} |w_j|^{\frac{1}{\gamma}} \log\Big( \frac{u_j}{u_{j-1}} \Big)\\
        &\le \|w_j\|_{\frac{1}{\gamma}}^{\frac{1}{\gamma}} \log\Big(\frac{1}{s'} \Big)\,.
    \end{align*}
    By assumption, $\gamma \le 1/2$ and so 
    \[
        \|w_j\|_{\frac{1}{\gamma}} \le \|w_j\|_{r'} \le 1.
    \]
    Combining the above with the estimate $\log(x) \le 1 + x$ for $x \ge 1$ bounds \eqref{e:I11} by 
    \begin{equation}
        \label{e:zeta_ub}
        2 \delta^{-\gamma} \Big(\frac{1 - s'}{s'}\Big)^\gamma \le 2 \delta^{-\alpha} \Big(\frac{1 - s'}{s'}\Big)^\gamma.
    \end{equation}
    The same argument applies to \eqref{e:I12}.
    It remains to deal with the $\mathcal{I}_2$ term \eqref{e:I2}. We  notice that for $j \in \mathcal{I}_2$ 
    \begin{equation}
        \label{e:tel}
        \big|\mathbf{1}_{[u_{j-1}, u_j]}(t) - \mathbf{1}_{[u_{j-1}, u_j]}(\frac{1}{s'}t)\big| = \mathbf{1}_{[s'u_{j-1}, u_{j-1}]}(t) + \mathbf{1}_{[s'u_j, u_j]}(t).
    \end{equation}
    Using this in \eqref{e:I2}, and subsequently doing the same Jensen's inequality and computation as for term \eqref{e:I11} also estimates \eqref{e:I2} by \eqref{e:zeta_ub}. This completes the proof. 
    \end{proof}

\end{document}
